Ecuador experienced one of the fastest escalations of lethal violence in contemporary Latin America: intentional homicides rose from 996 in 2018 to 9,283 in 2025 (a rate of approximately 51 per 100,000 inhabitants), with 4,154 additional deaths recorded in the first half of 2026. This paper examines the association between provincial labor market conditions and homicide rates using exclusively official sources. We build a province--quarter panel of 24 provinces over 34 quarters (816 observations, 2018Q1--2026Q2) from DINASED homicide records, INEC population projections, and the provincial cross-section of the ENEMDU labor force survey (2024Q2). Cross-sectional correlations between labor indicators and homicides are weak (Pearson r $\leq$ 0.34; Spearman $\rho$ = 0.42 for underemployment in 2025), and regression estimates with regional controls and clustered standard errors are small and statistically insignificant (unemployment $\beta$ = +0.08, p $\approx$ 0.39; underemployment $\beta$ = $-$0.04, p $\approx$ 0.47). An event study around the January 2024 "internal armed conflict" decree shows no significant differential change in homicides, and a placebo event in 2021Q1 is flat. A synthetic control for Esmeraldas indicates the province remained below its counterfactual after 2024 (mean gap $-$4.87 per 100,000), though placebo-based inference is coarse. The results are consistent with the international evidence that aggregate labor market deterioration does not drive homicide waves, which in Ecuador point to organized crime, illicit markets, and state capacity. Causal claims are explicitly out of reach with current data.